\documentclass[10pt,a4paper]{amsart}
\usepackage{amsmath,amssymb,amsthm,graphicx}
% For commutative diagrams we use Xy-pic
\usepackage[all]{xy}

% We use 2cell for 2-commutative diagrams.
\xyoption{2cell}
% Theorem environments.
%
\theoremstyle{plain}
\newtheorem{theorem}[subsection]{Theorem}
\newtheorem{proposition}[subsection]{Proposition}
\newtheorem{lemma}[subsection]{Lemma}

\theoremstyle{definition}
\newtheorem{definition}[subsection]{Definition}
\newtheorem{example}[subsection]{Example}
\newtheorem{exercise}[subsection]{Exercise}
\newtheorem{situation}[subsection]{Situation}

\theoremstyle{remark}
\newtheorem{remark}[subsection]{Remark}
\newtheorem{remarks}[subsection]{Remarks}

\numberwithin{equation}{subsection}

% Macros
%
\def\lim{\mathop{\mathrm{lim}}\nolimits}
\def\colim{\mathop{\mathrm{colim}}\nolimits}
\def\Spec{\mathop{\mathrm{Spec}}}
\def\Hom{\mathop{\mathrm{Hom}}\nolimits}
\def\Ext{\mathop{\mathrm{Ext}}\nolimits}
\def\SheafHom{\mathop{\mathcal{H}\!\mathit{om}}\nolimits}
\def\SheafExt{\mathop{\mathcal{E}\!\mathit{xt}}\nolimits}
\def\Sch{\mathit{Sch}}
\def\Mor{\mathop{\mathrm{Mor}}\nolimits}
\def\Ob{\mathop{\mathrm{Ob}}\nolimits}
\def\Sh{\mathop{\mathit{Sh}}\nolimits}
\def\NL{\mathop{N\!L}\nolimits}
\def\CH{\mathop{\mathrm{CH}}\nolimits}
\def\proetale{{pro\text{-}\acute{e}tale}}
\def\etale{{\acute{e}tale}}
\def\QCoh{\mathit{QCoh}}
\def\Ker{\mathop{\mathrm{Ker}}}
\def\Im{\mathop{\mathrm{Im}}}
\def\Coker{\mathop{\mathrm{Coker}}}
\def\Coim{\mathop{\mathrm{Coim}}}

% Boxtimes
%
\DeclareMathSymbol{\boxtimes}{\mathbin}{AMSa}{"02}

%
% Macros for moduli stacks/spaces
%
\def\QCohstack{\mathcal{QC}\!\mathit{oh}}
\def\Cohstack{\mathcal{C}\!\mathit{oh}}
\def\Spacesstack{\mathcal{S}\!\mathit{paces}}
\def\Quotfunctor{\mathrm{Quot}}
\def\Hilbfunctor{\mathrm{Hilb}}
\def\Curvesstack{\mathcal{C}\!\mathit{urves}}
\def\Polarizedstack{\mathcal{P}\!\mathit{olarized}}
\def\Complexesstack{\mathcal{C}\!\mathit{omplexes}}
% \Pic is the operator that assigns to X its picard group, usage \Pic(X)
% \Picardstack_{X/B} denotes the Picard stack of X over B
% \Picardfunctor_{X/B} denotes the Picard functor of X over B
\def\Pic{\mathop{\mathrm{Pic}}\nolimits}
\def\Picardstack{\mathcal{P}\!\mathit{ic}}
\def\Picardfunctor{\mathrm{Pic}}
\def\Deformationcategory{\mathcal{D}\!\mathit{ef}}
\newenvironment{reference}{}{}
\newenvironment{slogan}{}{}
\newenvironment{history}{}{}
\usepackage{cleveref}
\crefname{originaldefinition}{original definition}{original definitions}
\Crefname{originaldefinition}{Original definition}{Original definitions}
\makeatletter
\newlabel{definition-homology}{{01GB}{source}{}{sourcepointer-1}{}}
\newlabel{definition-homology@cref}{{[originaldefinition][0][]01GB}{[1][0][]source}{}{}{}}
\makeatother
\crefname{originallemma}{original lemma}{original lemmas}
\Crefname{originallemma}{Original lemma}{Original lemmas}
\makeatletter
\newlabel{lemma-hypercovering-acyclic}{{01GF}{source}{}{sourcepointer-2}{}}
\newlabel{lemma-hypercovering-acyclic@cref}{{[originallemma][0][]01GF}{[1][0][]source}{}{}{}}
\makeatother
\crefname{originallemma}{original lemma}{original lemmas}
\Crefname{originallemma}{Original lemma}{Original lemmas}
\makeatletter
\newlabel{sites-cohomology-lemma-cech-spectral-sequence}{{03AZ}{source}{}{sourcepointer-3}{}}
\newlabel{sites-cohomology-lemma-cech-spectral-sequence@cref}{{[originallemma][0][]03AZ}{[1][0][]source}{}{}{}}
\makeatother
\crefname{originalsection}{original section}{original sections}
\Crefname{originalsection}{Original section}{Original sections}
\makeatletter
\newlabel{homology-section-double-complexes}{{0FNB}{source}{}{sourcepointer-4}{}}
\newlabel{homology-section-double-complexes@cref}{{[originalsection][0][]0FNB}{[1][0][]source}{}{}{}}
\makeatother
\crefname{originalsection}{original section}{original sections}
\Crefname{originalsection}{Original section}{Original sections}
\makeatletter
\newlabel{homology-section-double-complex}{{012X}{source}{}{sourcepointer-5}{}}
\newlabel{homology-section-double-complex@cref}{{[originalsection][0][]012X}{[1][0][]source}{}{}{}}
\makeatother
\crefname{originallemma}{original lemma}{original lemmas}
\Crefname{originallemma}{Original lemma}{Original lemmas}
\makeatletter
\newlabel{homology-lemma-double-complex-gives-resolution}{{0133}{source}{}{sourcepointer-6}{}}
\newlabel{homology-lemma-double-complex-gives-resolution@cref}{{[originallemma][0][]0133}{[1][0][]source}{}{}{}}
\makeatother
\crefname{originallemma}{original lemma}{original lemmas}
\Crefname{originallemma}{Original lemma}{Original lemmas}
\makeatletter
\newlabel{homology-lemma-ss-double-complex}{{0130}{source}{}{sourcepointer-7}{}}
\newlabel{homology-lemma-ss-double-complex@cref}{{[originallemma][0][]0130}{[1][0][]source}{}{}{}}
\makeatother
\crefname{originallemma}{original lemma}{original lemmas}
\Crefname{originallemma}{Original lemma}{Original lemmas}
\makeatletter
\newlabel{homology-lemma-first-quadrant-ss}{{0132}{source}{}{sourcepointer-8}{}}
\newlabel{homology-lemma-first-quadrant-ss@cref}{{[originallemma][0][]0132}{[1][0][]source}{}{}{}}
\makeatother
\crefname{originalsection}{original section}{original sections}
\Crefname{originalsection}{Original section}{Original sections}
\makeatletter
\newlabel{simplicial-example-simplex-cosimplicial-set}{{0B13}{source}{}{sourcepointer-9}{}}
\newlabel{simplicial-example-simplex-cosimplicial-set@cref}{{[originalsection][0][]0B13}{[1][0][]source}{}{}{}}
\makeatother
\crefname{originallemma}{original lemma}{original lemmas}
\Crefname{originallemma}{Original lemma}{Original lemmas}
\makeatletter
\newlabel{simplicial-lemma-morphism-into-product}{{017M}{source}{}{sourcepointer-10}{}}
\newlabel{simplicial-lemma-morphism-into-product@cref}{{[originallemma][0][]017M}{[1][0][]source}{}{}{}}
\makeatother
\crefname{originallemma}{original lemma}{original lemmas}
\Crefname{originallemma}{Original lemma}{Original lemmas}
\makeatletter
\newlabel{lemma-funny-gamma}{{01GH}{source}{}{sourcepointer-11}{}}
\newlabel{lemma-funny-gamma@cref}{{[originallemma][0][]01GH}{[1][0][]source}{}{}{}}
\makeatother
\crefname{originalsection}{original section}{original sections}
\Crefname{originalsection}{Original section}{Original sections}
\makeatletter
\newlabel{simplicial-example-cosk0}{{0182}{source}{}{sourcepointer-12}{}}
\newlabel{simplicial-example-cosk0@cref}{{[originalsection][0][]0182}{[1][0][]source}{}{}{}}
\makeatother
\crefname{originallemma}{original lemma}{original lemmas}
\Crefname{originallemma}{Original lemma}{Original lemmas}
\makeatletter
\newlabel{simplicial-lemma-formula-limit}{{0186}{source}{}{sourcepointer-13}{}}
\newlabel{simplicial-lemma-formula-limit@cref}{{[originallemma][0][]0186}{[1][0][]source}{}{}{}}
\makeatother
\crefname{originallemma}{original lemma}{original lemmas}
\Crefname{originallemma}{Original lemma}{Original lemmas}
\makeatletter
\newlabel{simplicial-lemma-relations-face-degeneracy}{{0167}{source}{}{sourcepointer-14}{}}
\newlabel{simplicial-lemma-relations-face-degeneracy@cref}{{[originallemma][0][]0167}{[1][0][]source}{}{}{}}
\makeatother
\crefname{originalsection}{original section}{original sections}
\Crefname{originalsection}{Original section}{Original sections}
\makeatletter
\newlabel{example-cech}{{01G6}{source}{}{sourcepointer-15}{}}
\newlabel{example-cech@cref}{{[originalsection][0][]01G6}{[1][0][]source}{}{}{}}
\makeatother
\crefname{originallemma}{original lemma}{original lemmas}
\Crefname{originallemma}{Original lemma}{Original lemmas}
\makeatletter
\newlabel{sites-cohomology-lemma-kill-cohomology-class-on-covering}{{01FW}{source}{}{sourcepointer-16}{}}
\newlabel{sites-cohomology-lemma-kill-cohomology-class-on-covering@cref}{{[originallemma][0][]01FW}{[1][0][]source}{}{}{}}
\makeatother
\crefname{originalsection}{original section}{original sections}
\Crefname{originalsection}{Original section}{Original sections}
\makeatletter
\newlabel{sites-section-limits-colimits-PSh}{{00VB}{source}{}{sourcepointer-17}{}}
\newlabel{sites-section-limits-colimits-PSh@cref}{{[originalsection][0][]00VB}{[1][0][]source}{}{}{}}
\makeatother
\crefname{originallemma}{original lemma}{original lemmas}
\Crefname{originallemma}{Original lemma}{Original lemmas}
\makeatletter
\newlabel{sites-lemma-essential-image-j-shriek}{{00Y1}{source}{}{sourcepointer-18}{}}
\newlabel{sites-lemma-essential-image-j-shriek@cref}{{[originallemma][0][]00Y1}{[1][0][]source}{}{}{}}
\makeatother
\crefname{originallemma}{original lemma}{original lemmas}
\Crefname{originallemma}{Original lemma}{Original lemmas}
\makeatletter
\newlabel{categories-lemma-finite-limits-exist}{{002O}{source}{}{sourcepointer-19}{}}
\newlabel{categories-lemma-finite-limits-exist@cref}{{[originallemma][0][]002O}{[1][0][]source}{}{}{}}
\makeatother
\crefname{originallemma}{original lemma}{original lemmas}
\Crefname{originallemma}{Original lemma}{Original lemmas}
\makeatletter
\newlabel{simplicial-lemma-existence-cosk}{{0183}{source}{}{sourcepointer-20}{}}
\newlabel{simplicial-lemma-existence-cosk@cref}{{[originallemma][0][]0183}{[1][0][]source}{}{}{}}
\makeatother
\crefname{originalsection}{original section}{original sections}
\Crefname{originalsection}{Original section}{Original sections}
\makeatletter
\newlabel{simplicial-section-dold-kan-cosimplicial}{{019H}{source}{}{sourcepointer-21}{}}
\newlabel{simplicial-section-dold-kan-cosimplicial@cref}{{[originalsection][0][]019H}{[1][0][]source}{}{}{}}
\makeatother
\crefname{originalsection}{original section}{original sections}
\Crefname{originalsection}{Original section}{Original sections}
\makeatletter
\newlabel{cohomology-section-cech}{{01ED}{source}{}{sourcepointer-22}{}}
\newlabel{cohomology-section-cech@cref}{{[originalsection][0][]01ED}{[1][0][]source}{}{}{}}
\makeatother
\crefname{originalsection}{original section}{original sections}
\Crefname{originalsection}{Original section}{Original sections}
\makeatletter
\newlabel{cohomology-section-cech-functor}{{01EH}{source}{}{sourcepointer-23}{}}
\newlabel{cohomology-section-cech-functor@cref}{{[originalsection][0][]01EH}{[1][0][]source}{}{}{}}
\makeatother
\crefname{originalsection}{original section}{original sections}
\Crefname{originalsection}{Original section}{Original sections}
\makeatletter
\newlabel{cohomology-section-cech-cohomology-cohomology}{{01EO}{source}{}{sourcepointer-24}{}}
\newlabel{cohomology-section-cech-cohomology-cohomology@cref}{{[originalsection][0][]01EO}{[1][0][]source}{}{}{}}
\makeatother
\crefname{originallemma}{original lemma}{original lemmas}
\Crefname{originallemma}{Original lemma}{Original lemmas}
\makeatletter
\newlabel{simplicial-lemma-functorial-homotopy}{{019Y}{source}{}{sourcepointer-25}{}}
\newlabel{simplicial-lemma-functorial-homotopy@cref}{{[originallemma][0][]019Y}{[1][0][]source}{}{}{}}
\makeatother
\crefname{originallemma}{original lemma}{original lemmas}
\Crefname{originallemma}{Original lemma}{Original lemmas}
\makeatletter
\newlabel{simplicial-lemma-homotopy-s-Q}{{01A0}{source}{}{sourcepointer-26}{}}
\newlabel{simplicial-lemma-homotopy-s-Q@cref}{{[originallemma][0][]01A0}{[1][0][]source}{}{}{}}
\makeatother
\crefname{originallemma}{original lemma}{original lemmas}
\Crefname{originallemma}{Original lemma}{Original lemmas}
\makeatletter
\newlabel{lemma-hypercoverings-set}{{01G7}{source}{}{sourcepointer-27}{}}
\newlabel{lemma-hypercoverings-set@cref}{{[originallemma][0][]01G7}{[1][0][]source}{}{}{}}
\makeatother
\crefname{originalsection}{original section}{original sections}
\Crefname{originalsection}{Original section}{Original sections}
\makeatletter
\newlabel{categories-section-limits}{{002D}{source}{}{sourcepointer-28}{}}
\newlabel{categories-section-limits@cref}{{[originalsection][0][]002D}{[1][0][]source}{}{}{}}
\makeatother
\crefname{originallemma}{original lemma}{original lemmas}
\Crefname{originallemma}{Original lemma}{Original lemmas}
\makeatletter
\newlabel{lemma-product-hypercoverings}{{01GI}{source}{}{sourcepointer-29}{}}
\newlabel{lemma-product-hypercoverings@cref}{{[originallemma][0][]01GI}{[1][0][]source}{}{}{}}
\makeatother
\crefname{originallemma}{original lemma}{original lemmas}
\Crefname{originallemma}{Original lemma}{Original lemmas}
\makeatletter
\newlabel{lemma-homotopy}{{01GS}{source}{}{sourcepointer-30}{}}
\newlabel{lemma-homotopy@cref}{{[originallemma][0][]01GS}{[1][0][]source}{}{}{}}
\makeatother
\crefname{originaldefinition}{original definition}{original definitions}
\Crefname{originaldefinition}{Original definition}{Original definitions}
\makeatletter
\newlabel{categories-definition-directed}{{002V}{source}{}{sourcepointer-31}{}}
\newlabel{categories-definition-directed@cref}{{[originaldefinition][0][]002V}{[1][0][]source}{}{}{}}
\makeatother
\title{Sheaf cohomology is the filtered hypercover-Cech cohomology functor}
\author{The Stacks Project Authors}
\date{stacks-project:hypercovering.tex:local-snapshot-20261011; GNU FDL 1.2 or later}
\begin{document}
\maketitle
\noindent\textit{Editorial note.} This is a colimit over hypercover refinements. The separately prepared fixed-hypercover derived descent statement has a different target and requires an unbounded totalization argument; neither is counted as an alternate proof of the same assertion.\par\medskip
\noindent\textit{Allowed standard prerequisites (editorial).}\begin{itemize}
\item Explicit prerequisite boundary. Sites with fibre products, semi-representable objects, simplicial matching/coskeleton constructions, local annihilation of higher sheaf-cohomology classes and filtered colimits. The hypercover refinement andCech spectral-sequence proofs are retained.
\end{itemize}\medskip
\ifdefined\hypertarget\hypertarget{sourcepointer-1}{}\fi
\noindent\textit{Source reference.} Original definition 01GB: focused-seeds/AGP/sources/stacks-snapshot-20261011/tags.txt, source lines 1900--1900. Stable Stacks tag for an expressly allowed foundational or notational input; the counted construction remains the two concrete core obligations. The retained local core is not replaced by this pointer. See standard\_prerequisites.\par\medskip
\ifdefined\hypertarget\hypertarget{sourcepointer-2}{}\fi
\noindent\textit{Source reference.} Original lemma 01GF: focused-seeds/AGP/sources/stacks-snapshot-20261011/tags.txt, source lines 1904--1904. Stable Stacks tag for an expressly allowed foundational or notational input; the counted construction remains the two concrete core obligations. The retained local core is not replaced by this pointer. See standard\_prerequisites.\par\medskip
\ifdefined\hypertarget\hypertarget{sourcepointer-3}{}\fi
\noindent\textit{Source reference.} Original lemma 03AZ: focused-seeds/AGP/sources/stacks-snapshot-20261011/tags.txt, source lines 4328--4328. Stable Stacks tag for an expressly allowed foundational or notational input; the counted construction remains the two concrete core obligations. The retained local core is not replaced by this pointer. See standard\_prerequisites.\par\medskip
\ifdefined\hypertarget\hypertarget{sourcepointer-4}{}\fi
\noindent\textit{Source reference.} Original section 0FNB: focused-seeds/AGP/sources/stacks-snapshot-20261011/tags.txt, source lines 19561--19561. Stable Stacks tag for an expressly allowed foundational or notational input; the counted construction remains the two concrete core obligations. The retained local core is not replaced by this pointer. See standard\_prerequisites.\par\medskip
\ifdefined\hypertarget\hypertarget{sourcepointer-5}{}\fi
\noindent\textit{Source reference.} Original section 012X: focused-seeds/AGP/sources/stacks-snapshot-20261011/tags.txt, source lines 1415--1415. Stable Stacks tag for an expressly allowed foundational or notational input; the counted construction remains the two concrete core obligations. The retained local core is not replaced by this pointer. See standard\_prerequisites.\par\medskip
\ifdefined\hypertarget\hypertarget{sourcepointer-6}{}\fi
\noindent\textit{Source reference.} Original lemma 0133: focused-seeds/AGP/sources/stacks-snapshot-20261011/tags.txt, source lines 1421--1421. Stable Stacks tag for an expressly allowed foundational or notational input; the counted construction remains the two concrete core obligations. The retained local core is not replaced by this pointer. See standard\_prerequisites.\par\medskip
\ifdefined\hypertarget\hypertarget{sourcepointer-7}{}\fi
\noindent\textit{Source reference.} Original lemma 0130: focused-seeds/AGP/sources/stacks-snapshot-20261011/tags.txt, source lines 1418--1418. Stable Stacks tag for an expressly allowed foundational or notational input; the counted construction remains the two concrete core obligations. The retained local core is not replaced by this pointer. See standard\_prerequisites.\par\medskip
\ifdefined\hypertarget\hypertarget{sourcepointer-8}{}\fi
\noindent\textit{Source reference.} Original lemma 0132: focused-seeds/AGP/sources/stacks-snapshot-20261011/tags.txt, source lines 1420--1420. Stable Stacks tag for an expressly allowed foundational or notational input; the counted construction remains the two concrete core obligations. The retained local core is not replaced by this pointer. See standard\_prerequisites.\par\medskip
\ifdefined\hypertarget\hypertarget{sourcepointer-9}{}\fi
\noindent\textit{Source reference.} Original section 0B13: focused-seeds/AGP/sources/stacks-snapshot-20261011/tags.txt, source lines 13881--13881. Stable Stacks tag for an expressly allowed foundational or notational input; the counted construction remains the two concrete core obligations. The retained local core is not replaced by this pointer. See standard\_prerequisites.\par\medskip
\ifdefined\hypertarget\hypertarget{sourcepointer-10}{}\fi
\noindent\textit{Source reference.} Original lemma 017M: focused-seeds/AGP/sources/stacks-snapshot-20261011/tags.txt, source lines 1585--1585. Stable Stacks tag for an expressly allowed foundational or notational input; the counted construction remains the two concrete core obligations. The retained local core is not replaced by this pointer. See standard\_prerequisites.\par\medskip
\ifdefined\hypertarget\hypertarget{sourcepointer-11}{}\fi
\noindent\textit{Source reference.} Original lemma 01GH: focused-seeds/AGP/sources/stacks-snapshot-20261011/tags.txt, source lines 1906--1906. Stable Stacks tag for an expressly allowed foundational or notational input; the counted construction remains the two concrete core obligations. The retained local core is not replaced by this pointer. See standard\_prerequisites.\par\medskip
\ifdefined\hypertarget\hypertarget{sourcepointer-12}{}\fi
\noindent\textit{Source reference.} Original section 0182: focused-seeds/AGP/sources/stacks-snapshot-20261011/tags.txt, source lines 1601--1601. Stable Stacks tag for an expressly allowed foundational or notational input; the counted construction remains the two concrete core obligations. The retained local core is not replaced by this pointer. See standard\_prerequisites.\par\medskip
\ifdefined\hypertarget\hypertarget{sourcepointer-13}{}\fi
\noindent\textit{Source reference.} Original lemma 0186: focused-seeds/AGP/sources/stacks-snapshot-20261011/tags.txt, source lines 1605--1605. Stable Stacks tag for an expressly allowed foundational or notational input; the counted construction remains the two concrete core obligations. The retained local core is not replaced by this pointer. See standard\_prerequisites.\par\medskip
\ifdefined\hypertarget\hypertarget{sourcepointer-14}{}\fi
\noindent\textit{Source reference.} Original lemma 0167: focused-seeds/AGP/sources/stacks-snapshot-20261011/tags.txt, source lines 1534--1534. Stable Stacks tag for an expressly allowed foundational or notational input; the counted construction remains the two concrete core obligations. The retained local core is not replaced by this pointer. See standard\_prerequisites.\par\medskip
\ifdefined\hypertarget\hypertarget{sourcepointer-15}{}\fi
\noindent\textit{Source reference.} Original section 01G6: focused-seeds/AGP/sources/stacks-snapshot-20261011/tags.txt, source lines 1895--1895. Stable Stacks tag for an expressly allowed foundational or notational input; the counted construction remains the two concrete core obligations. The retained local core is not replaced by this pointer. See standard\_prerequisites.\par\medskip
\ifdefined\hypertarget\hypertarget{sourcepointer-16}{}\fi
\noindent\textit{Source reference.} Original lemma 01FW: focused-seeds/AGP/sources/stacks-snapshot-20261011/tags.txt, source lines 1885--1885. Stable Stacks tag for an expressly allowed foundational or notational input; the counted construction remains the two concrete core obligations. The retained local core is not replaced by this pointer. See standard\_prerequisites.\par\medskip
\ifdefined\hypertarget\hypertarget{sourcepointer-17}{}\fi
\noindent\textit{Source reference.} Original section 00VB: focused-seeds/AGP/sources/stacks-snapshot-20261011/tags.txt, source lines 1141--1141. Stable Stacks tag for an expressly allowed foundational or notational input; the counted construction remains the two concrete core obligations. The retained local core is not replaced by this pointer. See standard\_prerequisites.\par\medskip
\ifdefined\hypertarget\hypertarget{sourcepointer-18}{}\fi
\noindent\textit{Source reference.} Original lemma 00Y1: focused-seeds/AGP/sources/stacks-snapshot-20261011/tags.txt, source lines 1239--1239. Stable Stacks tag for an expressly allowed foundational or notational input; the counted construction remains the two concrete core obligations. The retained local core is not replaced by this pointer. See standard\_prerequisites.\par\medskip
\ifdefined\hypertarget\hypertarget{sourcepointer-19}{}\fi
\noindent\textit{Source reference.} Original lemma 002O: focused-seeds/AGP/sources/stacks-snapshot-20261011/tags.txt, source lines 105--105. Stable Stacks tag for an expressly allowed foundational or notational input; the counted construction remains the two concrete core obligations. The retained local core is not replaced by this pointer. See standard\_prerequisites.\par\medskip
\ifdefined\hypertarget\hypertarget{sourcepointer-20}{}\fi
\noindent\textit{Source reference.} Original lemma 0183: focused-seeds/AGP/sources/stacks-snapshot-20261011/tags.txt, source lines 1602--1602. Stable Stacks tag for an expressly allowed foundational or notational input; the counted construction remains the two concrete core obligations. The retained local core is not replaced by this pointer. See standard\_prerequisites.\par\medskip
\ifdefined\hypertarget\hypertarget{sourcepointer-21}{}\fi
\noindent\textit{Source reference.} Original section 019H: focused-seeds/AGP/sources/stacks-snapshot-20261011/tags.txt, source lines 1652--1652. Stable Stacks tag for an expressly allowed foundational or notational input; the counted construction remains the two concrete core obligations. The retained local core is not replaced by this pointer. See standard\_prerequisites.\par\medskip
\ifdefined\hypertarget\hypertarget{sourcepointer-22}{}\fi
\noindent\textit{Source reference.} Original section 01ED: focused-seeds/AGP/sources/stacks-snapshot-20261011/tags.txt, source lines 1830--1830. Stable Stacks tag for an expressly allowed foundational or notational input; the counted construction remains the two concrete core obligations. The retained local core is not replaced by this pointer. See standard\_prerequisites.\par\medskip
\ifdefined\hypertarget\hypertarget{sourcepointer-23}{}\fi
\noindent\textit{Source reference.} Original section 01EH: focused-seeds/AGP/sources/stacks-snapshot-20261011/tags.txt, source lines 1834--1834. Stable Stacks tag for an expressly allowed foundational or notational input; the counted construction remains the two concrete core obligations. The retained local core is not replaced by this pointer. See standard\_prerequisites.\par\medskip
\ifdefined\hypertarget\hypertarget{sourcepointer-24}{}\fi
\noindent\textit{Source reference.} Original section 01EO: focused-seeds/AGP/sources/stacks-snapshot-20261011/tags.txt, source lines 1841--1841. Stable Stacks tag for an expressly allowed foundational or notational input; the counted construction remains the two concrete core obligations. The retained local core is not replaced by this pointer. See standard\_prerequisites.\par\medskip
\ifdefined\hypertarget\hypertarget{sourcepointer-25}{}\fi
\noindent\textit{Source reference.} Original lemma 019Y: focused-seeds/AGP/sources/stacks-snapshot-20261011/tags.txt, source lines 1669--1669. Stable Stacks tag for an expressly allowed foundational or notational input; the counted construction remains the two concrete core obligations. The retained local core is not replaced by this pointer. See standard\_prerequisites.\par\medskip
\ifdefined\hypertarget\hypertarget{sourcepointer-26}{}\fi
\noindent\textit{Source reference.} Original lemma 01A0: focused-seeds/AGP/sources/stacks-snapshot-20261011/tags.txt, source lines 1671--1671. Stable Stacks tag for an expressly allowed foundational or notational input; the counted construction remains the two concrete core obligations. The retained local core is not replaced by this pointer. See standard\_prerequisites.\par\medskip
\ifdefined\hypertarget\hypertarget{sourcepointer-27}{}\fi
\noindent\textit{Source reference.} Original lemma 01G7: focused-seeds/AGP/sources/stacks-snapshot-20261011/tags.txt, source lines 1896--1896. Stable Stacks tag for an expressly allowed foundational or notational input; the counted construction remains the two concrete core obligations. The retained local core is not replaced by this pointer. See standard\_prerequisites.\par\medskip
\ifdefined\hypertarget\hypertarget{sourcepointer-28}{}\fi
\noindent\textit{Source reference.} Original section 002D: focused-seeds/AGP/sources/stacks-snapshot-20261011/tags.txt, source lines 94--94. Stable Stacks tag for an expressly allowed foundational or notational input; the counted construction remains the two concrete core obligations. The retained local core is not replaced by this pointer. See standard\_prerequisites.\par\medskip
\ifdefined\hypertarget\hypertarget{sourcepointer-29}{}\fi
\noindent\textit{Source reference.} Original lemma 01GI: focused-seeds/AGP/sources/stacks-snapshot-20261011/tags.txt, source lines 1907--1907. Stable Stacks tag for an expressly allowed foundational or notational input; the counted construction remains the two concrete core obligations. The retained local core is not replaced by this pointer. See standard\_prerequisites.\par\medskip
\ifdefined\hypertarget\hypertarget{sourcepointer-30}{}\fi
\noindent\textit{Source reference.} Original lemma 01GS: focused-seeds/AGP/sources/stacks-snapshot-20261011/tags.txt, source lines 1917--1917. Stable Stacks tag for an expressly allowed foundational or notational input; the counted construction remains the two concrete core obligations. The retained local core is not replaced by this pointer. See standard\_prerequisites.\par\medskip
\ifdefined\hypertarget\hypertarget{sourcepointer-31}{}\fi
\noindent\textit{Source reference.} Original definition 002V: focused-seeds/AGP/sources/stacks-snapshot-20261011/tags.txt, source lines 112--112. Stable Stacks tag for an expressly allowed foundational or notational input; the counted construction remains the two concrete core obligations. The retained local core is not replaced by this pointer. See standard\_prerequisites.\par\medskip
\begin{theorem}
\label{theorem-cohomology-hypercoverings}
Let $\mathcal{C}$ be a site with fibre products.
Let $X$ be an object of $\mathcal{C}$. Let $i \geq 0$.
The functors
\begin{eqnarray*}
\textit{Ab}(\mathcal{C}) & \longrightarrow & \textit{Ab} \\
\mathcal{F} & \longmapsto & H^i(X, \mathcal{F}) \\
\mathcal{F} & \longmapsto & \check{H}^i_{\text{HC}}(X, \mathcal{F})
\end{eqnarray*}
are canonically isomorphic.
\end{theorem}

\section{Semi-representable objects}
\label{section-semi-representable}

\noindent
In order to start we make the following definition.
The letters ``SR'' stand for Semi-Representable.

\begin{definition}
\label{definition-SR}
Let $\mathcal{C}$ be a category. We denote $\text{SR}(\mathcal{C})$
the category of {\it semi-representable objects} defined as follows
\begin{enumerate}
\item objects are families of objects $\{U_i\}_{i \in I}$, and
\item morphisms $\{U_i\}_{i \in I} \to \{V_j\}_{j \in J}$ are given by
a map $\alpha : I \to J$ and for each $i \in I$
a morphism $f_i : U_i \to V_{\alpha(i)}$ of $\mathcal{C}$.
\end{enumerate}
Let $X \in \Ob(\mathcal{C})$ be an object of $\mathcal{C}$.
The category of {\it semi-representable objects over $X$}
is the category
$\text{SR}(\mathcal{C}, X) = \text{SR}(\mathcal{C}/X)$.
\end{definition}

\noindent
This definition is essentially equivalent to
\cite[Expos\'e V, Subsection 7.3.0]{SGA4}. Note that
this is a ``big'' category. We will later ``bound'' the size of the index
sets $I$ that we need for hypercoverings of $X$. We can then redefine
$\text{SR}(\mathcal{C}, X)$ to become a category. Let's spell out
the objects and morphisms $\text{SR}(\mathcal{C}, X)$:
\begin{enumerate}
\item objects are families of morphisms
$\{U_i \to X\}_{i \in I}$, and
\item morphisms $\{U_i \to X\}_{i \in I} \to
\{V_j \to X\}_{j \in J}$ are given by
a map $\alpha : I \to J$ and for each $i \in I$
a morphism $f_i : U_i \to V_{\alpha(i)}$ over $X$.
\end{enumerate}
There is a forgetful functor
$\text{SR}(\mathcal{C}, X) \to \text{SR}(\mathcal{C})$.

\begin{definition}
\label{definition-SR-F}
Let $\mathcal{C}$ be a category.
We denote $F$ the functor {\it which associates a presheaf to a
semi-representable object}. In a formula
\begin{eqnarray*}
F : \text{SR}(\mathcal{C}) & \longrightarrow & \textit{PSh}(\mathcal{C}) \\
\{U_i\}_{i \in I} & \longmapsto & \amalg_{i\in I} h_{U_i}
\end{eqnarray*}
where $h_U$ denotes the representable presheaf associated to
the object $U$.
\end{definition}

\noindent
Given a morphism $U \to X$ we obtain a morphism $h_U \to h_X$ of representable
presheaves. Thus we often think of $F$ on $\text{SR}(\mathcal{C}, X)$
as a functor into the category of presheaves of sets over $h_X$,
namely $\textit{PSh}(\mathcal{C})/h_X$. Here is a picture:
$$
\xymatrix{
\text{SR}(\mathcal{C}, X) \ar[r]_F \ar[d] &
\textit{PSh}(\mathcal{C})/h_X \ar[d] \\
\text{SR}(\mathcal{C}) \ar[r]^F &
\textit{PSh}(\mathcal{C})
}
$$
Next we discuss the existence of limits in the category of semi-representable
objects.

\begin{lemma}
\label{lemma-coprod-prod-SR}
Let $\mathcal{C}$ be a category.
\begin{enumerate}
\item the category $\text{SR}(\mathcal{C})$ has coproducts
and $F$ commutes with them,
\item the functor $F : \text{SR}(\mathcal{C}) \to \textit{PSh}(\mathcal{C})$
commutes with limits,
\item if $\mathcal{C}$ has fibre products, then $\text{SR}(\mathcal{C})$
has fibre products,
\item if $\mathcal{C}$ has products of pairs, then
$\text{SR}(\mathcal{C})$ has products of pairs,
\item if $\mathcal{C}$ has equalizers, so does $\text{SR}(\mathcal{C})$, and
\item if $\mathcal{C}$ has a final object, so does $\text{SR}(\mathcal{C})$.
\end{enumerate}
Let $X \in \Ob(\mathcal{C})$.
\begin{enumerate}
\item the category $\text{SR}(\mathcal{C}, X)$ has coproducts
and $F$ commutes with them,
\item if $\mathcal{C}$ has fibre products, then $\text{SR}(\mathcal{C}, X)$
has finite limits and
$F : \text{SR}(\mathcal{C}, X) \to \textit{PSh}(\mathcal{C})/h_X$
commutes with them.
\end{enumerate}
\end{lemma}

\begin{proof}
Proof of the results on $\text{SR}(\mathcal{C})$.
Proof of (1). The coproduct of $\{U_i\}_{i \in I}$ and $\{V_j\}_{j \in J}$ is
$\{U_i\}_{i \in I} \amalg \{V_j\}_{j \in J}$, in other words, the family
of objects whose index set is $I \amalg J$ and for an element
$k \in I \amalg J$ gives $U_i$ if $k = i \in I$ and gives $V_j$ if
$k = j \in J$. Similarly for coproducts
of families of objects. It is clear that $F$ commutes with these.

\medskip\noindent
Proof of (2). For $U$ in $\Ob(\mathcal{C})$ consider the object $\{U\}$ of
$\text{SR}(\mathcal{C})$. It is clear that
$\Mor_{\text{SR}(\mathcal{C})}(\{U\}, K)) = F(K)(U)$
for $K \in \Ob(\text{SR}(\mathcal{C}))$. Since limits of presheaves
are computed at the level of sections
(Sites, Section \ref{sites-section-limits-colimits-PSh})
we conclude that $F$ commutes with limits.

\medskip\noindent
Proof of (3). Suppose given a morphism
$(\alpha, f_i) : \{U_i\}_{i \in I} \to \{V_j\}_{j \in J}$
and a morphism
$(\beta, g_k) : \{W_k\}_{k \in K} \to \{V_j\}_{j \in J}$.
The fibred product of these morphisms is given by
$$
\{ U_i \times_{f_i, V_j, g_k} W_k\}_{(i, j, k) \in I \times J \times K
\text{ such that } j = \alpha(i) = \beta(k)}
$$
The fibre products exist if $\mathcal{C}$ has fibre products.

\medskip\noindent
Proof of (4). The product of $\{U_i\}_{i \in I}$ and $\{V_j\}_{j \in J}$ is
$\{U_i \times V_j\}_{i \in I, j \in J}$. The products exist if
$\mathcal{C}$ has products.

\medskip\noindent
Proof of (5). The equalizer of two maps
$(\alpha, f_i), (\alpha', f'_i) : \{U_i\}_{i \in I} \to \{V_j\}_{j \in J}$
is
$$
\{
\text{Eq}(f_i, f'_i : U_i \to V_{\alpha(i)})
\}_{i \in I,\ \alpha(i) = \alpha'(i)}
$$
The equalizers exist if $\mathcal{C}$ has equalizers.

\medskip\noindent
Proof of (6). If $X$ is a final object of $\mathcal{C}$, then
$\{X\}$ is a final object of $\text{SR}(\mathcal{C})$.

\medskip\noindent
Proof of the statements about $\text{SR}(\mathcal{C}, X)$.
These follow from the results above applied to the category
$\mathcal{C}/X$ using that
$\text{SR}(\mathcal{C}/X) = \text{SR}(\mathcal{C}, X)$ and that
$\textit{PSh}(\mathcal{C}/X) = \textit{PSh}(\mathcal{C})/h_X$
(Sites, Lemma \ref{sites-lemma-essential-image-j-shriek} applied
to $\mathcal{C}$ endowed with the chaotic topology). However
we also argue directly as follows.
It is clear that the coproduct of
$\{U_i \to X\}_{i \in I}$ and $\{V_j \to X\}_{j \in J}$
is $\{U_i \to X\}_{i \in I} \amalg \{V_j \to X\}_{j \in J}$
and similarly for coproducts of
families of families of morphisms with target $X$.
The object $\{X \to X\}$ is a final
object of $\text{SR}(\mathcal{C}, X)$.
Suppose given a morphism
$(\alpha, f_i) : \{U_i \to X\}_{i \in I} \to \{V_j \to X\}_{j \in J}$
and a morphism
$(\beta, g_k) : \{W_k \to X\}_{k \in K} \to \{V_j \to X\}_{j \in J}$.
The fibred product of these morphisms is given by
$$
\{ U_i \times_{f_i, V_j, g_k} W_k \to X \}_{(i, j, k) \in I \times J \times K
\text{ such that } j = \alpha(i) = \beta(k)}
$$
The fibre products exist by the assumption that
$\mathcal{C}$ has fibre products.
Thus $\text{SR}(\mathcal{C}, X)$ has finite limits,
see Categories, Lemma \ref{categories-lemma-finite-limits-exist}.
We omit verifying the statements on the functor $F$ in this case.
\end{proof}





\section{Hypercoverings}
\label{section-hypercoverings}

\noindent
If we assume our category is a site, then we can make the following
definition.

\begin{definition}
\label{definition-covering-SR}
Let $\mathcal{C}$ be a site. Let
$f = (\alpha, f_i) : \{U_i\}_{i \in I} \to \{V_j\}_{j \in J}$
be a morphism in the category $\text{SR}(\mathcal{C})$.
We say that $f$ is a {\it covering} if for every $j \in J$ the
family of morphisms $\{U_i \to V_j\}_{i \in I, \alpha(i) = j}$
is a covering for the site $\mathcal{C}$.
Let $X$ be an object of $\mathcal{C}$.
A morphism $K \to L$ in $\text{SR}(\mathcal{C}, X)$ is
a {\it covering} if its image in $\text{SR}(\mathcal{C})$ is
a covering.
\end{definition}

\begin{lemma}
\label{lemma-covering-permanence}
Let $\mathcal{C}$ be a site.
\begin{enumerate}
\item A composition of coverings in $\text{SR}(\mathcal{C})$
is a covering.
\item If $K \to L$ is a covering in $\text{SR}(\mathcal{C})$
and $L' \to L$ is a morphism, then $L' \times_L K$ exists
and $L' \times_L K \to L'$ is a covering.
\item If $\mathcal{C}$ has products of pairs, and
$A \to B$ and $K \to L$ are coverings in $\text{SR}(\mathcal{C})$,
then $A \times K \to B \times L$ is a covering.
\end{enumerate}
Let $X \in \Ob(\mathcal{C})$. Then (1) and (2) holds for
$\text{SR}(\mathcal{C}, X)$ and (3) holds if $\mathcal{C}$
has fibre products.
\end{lemma}

\begin{proof}
Part (1) is immediate from the axioms of a site.
Part (2) follows by the construction of fibre products
in $\text{SR}(\mathcal{C})$ in the proof of
Lemma \ref{lemma-coprod-prod-SR}
and the requirement that the morphisms in a covering
of $\mathcal{C}$ are representable.
Part (3) follows by thinking of $A \times K \to B \times L$
as the composition $A \times K \to B \times K \to B \times L$
and hence a composition of basechanges of coverings.
The final statement follows because $\text{SR}(\mathcal{C}, X) =
\text{SR}(\mathcal{C}/X)$.
\end{proof}

\noindent
By Lemma \ref{lemma-coprod-prod-SR} and
Simplicial, Lemma \ref{simplicial-lemma-existence-cosk}
the coskeleton of a truncated simplicial object of
$\text{SR}(\mathcal{C}, X)$ exists if $\mathcal{C}$ has fibre products.
Hence the following definition makes sense.

\begin{definition}
\label{definition-hypercovering}
Let $\mathcal{C}$ be a site. Assume $\mathcal{C}$ has fibre products.
Let $X \in \Ob(\mathcal{C})$ be an object of $\mathcal{C}$.
A {\it hypercovering of $X$} is a simplicial object
$K$ of $\text{SR}(\mathcal{C}, X)$ such that
\begin{enumerate}
\item The object $K_0$ is a covering of $X$ for the site $\mathcal{C}$.
\item For every $n \geq 0$ the canonical morphism
$$
K_{n + 1} \longrightarrow (\text{cosk}_n \text{sk}_n K)_{n + 1}
$$
is a covering in the sense defined above.
\end{enumerate}
\end{definition}

\section{{\v C}ech cohomology and hypercoverings}
\label{section-hyper-cech}

\noindent
Let $\mathcal{C}$ be a site. Consider a presheaf of
abelian groups $\mathcal{F}$ on the site $\mathcal{C}$.
It defines a functor
\begin{eqnarray*}
\mathcal{F} : \text{SR}(\mathcal{C})^{opp}
& \longrightarrow &
\textit{Ab} \\
\{U_i\}_{i \in I} &
\longmapsto &
\prod\nolimits_{i \in I} \mathcal{F}(U_i)
\end{eqnarray*}
Thus a simplicial object $K$ of $\text{SR}(\mathcal{C})$
is turned into a cosimplicial object $\mathcal{F}(K)$ of $\textit{Ab}$.
The cochain complex $s(\mathcal{F}(K))$ associated to $\mathcal{F}(K)$
(Simplicial, Section
\ref{simplicial-section-dold-kan-cosimplicial})
is called the {\v C}ech complex of $\mathcal{F}$ with
respect to the simplicial object $K$. We set
$$
\check{H}^i(K, \mathcal{F})
=
H^i(s(\mathcal{F}(K))).
$$
and we call it the $i$th {\v C}ech cohomology group
of $\mathcal{F}$ with respect to $K$.
In this section we prove analogues of some of the results for
{\v C}ech cohomology of open coverings proved in
Cohomology, Sections \ref{cohomology-section-cech},
\ref{cohomology-section-cech-functor} and
\ref{cohomology-section-cech-cohomology-cohomology}.


\section{Cohomology and hypercoverings}
\label{section-cohomology}

\noindent
Let $\mathcal{C}$ be a site with fibre products.
Let $X$ be an object of $\mathcal{C}$.
Let $\mathcal{F}$ be a sheaf of abelian groups on $\mathcal{C}$.
Let $K, L$ be hypercoverings of $X$.
If $a, b : K \to L$ are homotopic maps,
then $\mathcal{F}(a), \mathcal{F}(b) : \mathcal{F}(K) \to \mathcal{F}(L)$
are homotopic maps, see
Simplicial, Lemma \ref{simplicial-lemma-functorial-homotopy}.
Hence have the same effect on cohomology groups of the associated
cochain complexes, see
Simplicial, Lemma \ref{simplicial-lemma-homotopy-s-Q}.
We are going to use this to define the colimit over all
hypercoverings.

\medskip\noindent
Let us temporarily denote $\text{HC}(\mathcal{C}, X)$
the category whose objects are hypercoverings of $X$
and whose morphisms are maps between hypercoverings of $X$
up to homotopy. We have seen that
this is a category and not a ``big'' category,
see Lemma \ref{lemma-hypercoverings-set}.
The opposite to $\text{HC}(\mathcal{C}, X)$
will be the index category for our diagram, see
Categories, Section \ref{categories-section-limits} for terminology.
Consider the diagram
$$
\check{H}^i(-, \mathcal{F}) :
\text{HC}(\mathcal{C}, X)^{opp}
\longrightarrow
\textit{Ab}.
$$
By Lemmas \ref{lemma-product-hypercoverings} and \ref{lemma-homotopy}
and the remark on homotopies above, this diagram is directed, see
Categories, Definition \ref{categories-definition-directed}.
Thus the colimit
$$
\check{H}^i_{\text{HC}}(X, \mathcal{F}) =
\colim_{K \in \text{HC}(\mathcal{C}, X)} \check{H}^i(K, \mathcal{F})
$$
has a particularly simple description (see location cited).

\begin{lemma}
\label{lemma-h0-cech}
Let $\mathcal{C}$ be a site with fibre products.
Let $X$ be an object of $\mathcal{C}$.
Let $K$ be a hypercovering of $X$.
Let $\mathcal{F}$ be a sheaf of abelian groups on $\mathcal{C}$.
Then $\check{H}^0(K, \mathcal{F}) = \mathcal{F}(X)$.
\end{lemma}

\begin{proof}
We have
$$
\check{H}^0(K, \mathcal{F})
=
\Ker(\mathcal{F}(K_0) \longrightarrow \mathcal{F}(K_1))
$$
Write $K_0 = \{U_i \to X\}$. It is a covering in the site
$\mathcal{C}$. As well, we have that $K_1 \to K_0 \times K_0$
is a covering in $\text{SR}(\mathcal{C}, X)$. Hence we may
write $K_1 = \amalg_{i_0, i_1 \in I} \{V_{i_0i_1j} \to X\}$
so that the morphism $K_1 \to K_0 \times K_0$ is given
by coverings $\{V_{i_0i_1j} \to U_{i_0} \times_X U_{i_1}\}$
of the site $\mathcal{C}$. Thus we can further identify
$$
\check{H}^0(K, \mathcal{F})
=
\Ker(
\prod\nolimits_i \mathcal{F}(U_i)
\longrightarrow
\prod\nolimits_{i_0i_1 j} \mathcal{F}(V_{i_0i_1j})
)
$$
with obvious map. The sheaf property of $\mathcal{F}$
implies that $\check{H}^0(K, \mathcal{F}) = H^0(X, \mathcal{F})$.
\end{proof}

\noindent
In fact this property characterizes the abelian sheaves among all
abelian presheaves on $\mathcal{C}$ of course.
The analogue of Cohomology, Lemma \ref{lemma-injective-trivial-cech}
in this case is the following.

\begin{lemma}
\label{lemma-injective-trivial-cech}
Let $\mathcal{C}$ be a site with fibre products.
Let $X$ be an object of $\mathcal{C}$.
Let $K$ be a hypercovering of $X$.
Let $\mathcal{I}$ be an injective sheaf of abelian groups on $\mathcal{C}$.
Then
$$
\check{H}^p(K, \mathcal{I}) =
\left\{
\begin{matrix}
\mathcal{I}(X) & \text{if} & p = 0 \\
0 & \text{if} & p > 0
\end{matrix}
\right.
$$
\end{lemma}

\begin{proof}
Observe that for any object $Z = \{U_i \to X\}$ of
$\text{SR}(\mathcal{C}, X)$ and any abelian sheaf
$\mathcal{F}$ on $\mathcal{C}$ we have
\begin{eqnarray*}
\mathcal{F}(Z)
& = &
\prod \mathcal{F}(U_i) \\
& = &
\prod \Mor_{\textit{PSh}(\mathcal{C})}(h_{U_i}, \mathcal{F})\\
& = &
\Mor_{\textit{PSh}(\mathcal{C})}(F(Z), \mathcal{F})\\
& = &
\Mor_{\textit{PAb}(\mathcal{C})}(\mathbf{Z}_{F(Z)}, \mathcal{F}) \\
& = &
\Mor_{\textit{Ab}(\mathcal{C})}(\mathbf{Z}_{F(Z)}^\#, \mathcal{F})
\end{eqnarray*}
Thus we see, for any simplicial object $K$ of
$\text{SR}(\mathcal{C}, X)$ that we have
\begin{equation}
\label{equation-identify-cech}
s(\mathcal{F}(K))
=
\Hom_{\textit{Ab}(\mathcal{C})}(s(\mathbf{Z}_{F(K)}^\#), \mathcal{F})
\end{equation}
see Definition \ref{definition-homology} for notation.
The complex of sheaves $s(\mathbf{Z}_{F(K)}^\#)$ is quasi-isomorphic
to $\mathbf{Z}_X^\#$ if $K$ is a hypercovering, see
Lemma \ref{lemma-hypercovering-acyclic}. We conclude
that if $\mathcal{I}$ is an injective abelian sheaf, and
$K$ a hypercovering, then the complex $s(\mathcal{I}(K))$
is acyclic except possibly in degree $0$.
In other words, we have
$$
\check{H}^i(K, \mathcal{I}) = 0
$$
for $i > 0$. Combined with Lemma \ref{lemma-h0-cech} the lemma is proved.
\end{proof}

\noindent
Next we come to the analogue of Cohomology on Sites, Lemma
\ref{sites-cohomology-lemma-cech-spectral-sequence}.
Let $\mathcal{C}$ be a site.
Let $\mathcal{F}$ be a sheaf of abelian groups on $\mathcal{C}$.
Recall that $\underline{H}^i(\mathcal{F})$ indicates the presheaf
of abelian groups on $\mathcal{C}$ which is defined by the
rule $\underline{H}^i(\mathcal{F}) : U \longmapsto H^i(U, \mathcal{F})$.
We extend this to $\text{SR}(\mathcal{C})$ as in the introduction
to this section.

\begin{lemma}
\label{lemma-cech-spectral-sequence}
Let $\mathcal{C}$ be a site with fibre products.
Let $X$ be an object of $\mathcal{C}$.
Let $K$ be a hypercovering of $X$.
Let $\mathcal{F}$ be a sheaf of abelian groups on $\mathcal{C}$.
There is a map
$$
s(\mathcal{F}(K))
\longrightarrow
R\Gamma(X, \mathcal{F})
$$
in $D^{+}(\textit{Ab})$ functorial in $\mathcal{F}$, which induces
natural transformations
$$
\check{H}^i(K, -) \longrightarrow H^i(X, -)
$$
as functors $\textit{Ab}(\mathcal{C}) \to \textit{Ab}$. Moreover,
there is a spectral sequence $(E_r, d_r)_{r \geq 0}$ with
$$
E_2^{p, q} = \check{H}^p(K, \underline{H}^q(\mathcal{F}))
$$
converging to $H^{p + q}(X, \mathcal{F})$.
This spectral sequence is functorial in $\mathcal{F}$ and
in the hypercovering $K$.
\end{lemma}

\begin{proof}
We could prove this by the same method as employed in the corresponding
lemma in the chapter on cohomology. Instead let us prove this by a
double complex argument.

\medskip\noindent
Choose an injective resolution $\mathcal{F} \to \mathcal{I}^\bullet$
in the category of abelian sheaves on $\mathcal{C}$. Consider the
double complex $A^{\bullet, \bullet}$ with terms
$$
A^{p, q} = \mathcal{I}^q(K_p)
$$
where the differential $d_1^{p, q} : A^{p, q} \to A^{p + 1, q}$ is the one
coming from the differential on the complex $s(\mathcal{I}^q(K))$
associated to the cosimplicial abelian group $\mathcal{I}^p(K)$
and the differential $d_2^{p, q} : A^{p, q} \to A^{p, q + 1}$ is the one
coming from the differential $\mathcal{I}^q \to \mathcal{I}^{q + 1}$.
Denote $\text{Tot}(A^{\bullet, \bullet})$ the total complex associated to
the double complex $A^{\bullet, \bullet}$, see
Homology, Section \ref{homology-section-double-complexes}.
We will use the two spectral
sequences $({}'E_r, {}'d_r)$ and $({}''E_r, {}''d_r)$
associated to this double complex, see
Homology, Section \ref{homology-section-double-complex}.

\medskip\noindent
By Lemma \ref{lemma-injective-trivial-cech}
the complexes $s(\mathcal{I}^q(K))$ are acyclic in
positive degrees and have $H^0$ equal to $\mathcal{I}^q(X)$.
Hence by
Homology, Lemma \ref{homology-lemma-double-complex-gives-resolution}
the natural map
$$
\mathcal{I}^\bullet(X) \longrightarrow \text{Tot}(A^{\bullet, \bullet})
$$
is a quasi-isomorphism of complexes of abelian groups. In particular
we conclude that $H^n(\text{Tot}(A^{\bullet, \bullet})) = H^n(X, \mathcal{F})$.

\medskip\noindent
The map $s(\mathcal{F}(K)) \longrightarrow R\Gamma(X, \mathcal{F})$ of
the lemma is the composition of the map
$s(\mathcal{F}(K)) \to \text{Tot}(A^{\bullet, \bullet})$
followed by the inverse
of the displayed quasi-isomorphism above. This works because
$\mathcal{I}^\bullet(X)$ is a representative of $R\Gamma(X, \mathcal{F})$.

\medskip\noindent
Consider the spectral sequence $({}'E_r, {}'d_r)_{r \geq 0}$. By
Homology, Lemma \ref{homology-lemma-ss-double-complex}
we see that
$$
{}'E_2^{p, q} = H^p_I(H^q_{II}(A^{\bullet, \bullet}))
$$
In other words, we first take cohomology with respect to
$d_2$ which gives the groups
${}'E_1^{p, q} = \underline{H}^q(\mathcal{F})(K_p)$.
Hence it is indeed the case (by the description of the differential
${}'d_1$) that
${}'E_2^{p, q} = \check{H}^p(K, \underline{H}^q(\mathcal{F}))$.
By the above and Homology, Lemma \ref{homology-lemma-first-quadrant-ss}
we see that this converges to $H^n(X, \mathcal{F})$ as desired.

\medskip\noindent
We omit the proof of the statements regarding the functoriality of
the above constructions in the abelian sheaf $\mathcal{F}$ and the
hypercovering $K$.
\end{proof}

\begin{lemma}
\label{lemma-covering}
Let $\mathcal{C}$ be a site with fibre products.
Let $X$ be an object of $\mathcal{C}$.
Let $K$ be a hypercovering of $X$.
Let $k \geq 0$ be an integer.
Let $u : Z \to K_k$ be a covering
in $\text{SR}(\mathcal{C}, X)$.
Then there exists a morphism of hypercoverings
$f: L \to K$ such that $L_k \to K_k$
factors through $u$.
\end{lemma}

\begin{proof}
Denote $Y = K_k$. Let $C[k]$ be the cosimplicial set defined in
Simplicial, Example \ref{simplicial-example-simplex-cosimplicial-set}.
We will use the description of $\Hom(C[k], Y)$ and $\Hom(C[k], Z)$
given in
Simplicial, Lemma \ref{simplicial-lemma-morphism-into-product}.
There is a canonical morphism
$K \to \Hom(C[k], Y)$ corresponding to $\text{id} : K_k = Y \to Y$.
Consider the morphism $\Hom(C[k], Z) \to \Hom(C[k], Y)$
which on degree $n$ terms is the morphism
$$
\prod\nolimits_{\alpha : [k] \to [n]} Z
\longrightarrow
\prod\nolimits_{\alpha : [k] \to [n]} Y
$$
using the given morphism $Z \to Y$ on each factor. Set
$$
L = K \times_{\Hom(C[k], Y)} \Hom(C[k], Z).
$$
The morphism $L_k \to K_k$ sits in to a commutative diagram
$$
\xymatrix{
L_k \ar[r] \ar[d] &
\prod_{\alpha : [k] \to [k]} Z \ar[r]^-{\text{pr}_{\text{id}_{[k]}}} \ar[d] &
Z \ar[d] \\
K_k \ar[r] &
\prod_{\alpha : [k] \to [k]} Y \ar[r]^-{\text{pr}_{\text{id}_{[k]}}} &
Y
}
$$
Since the composition of the two bottom arrows is the identity
we conclude that we have the desired factorization.

\medskip\noindent
We still have to show that $L$ is a hypercovering of $X$.
To see this we will use Lemma \ref{lemma-funny-gamma}.
Condition (1) is satisfied by assumption.
For (2), the morphism
$$
\Hom(C[k], Z)_0 \to \Hom(C[k], Y)_0
$$
is a covering because it is isomorphic to $Z \to Y$ as
there is only one morphism $[k] \to [0]$.

\medskip\noindent
Let us consider condition (3) for $n = 0$. Then, since
$(\text{cosk}_0 T)_1 = T \times T$
(Simplicial, Example \ref{simplicial-example-cosk0})
and since $\Hom(C[k], Z)_1 = \prod_{\alpha : [k] \to [1]} Z$
we obtain the diagram
$$
\xymatrix{
\prod\nolimits_{\alpha : [k] \to [1]} Z \ar[r] \ar[d] &
Z \times Z \ar[d] \\
\prod\nolimits_{\alpha : [k] \to [1]} Y \ar[r] &
Y \times Y
}
$$
with horizontal arrows corresponding to the projection onto the factors
corresponding to the two nonsurjective $\alpha$. Thus the arrow $\gamma$
is the morphism
$$
\prod\nolimits_{\alpha : [k] \to [1]} Z
\longrightarrow
\prod\nolimits_{\alpha : [k] \to [1]\text{ not onto}} Z
\times
\prod\nolimits_{\alpha : [k] \to [1]\text{ onto}} Y
$$
which is a product of coverings and hence a covering by
Lemma \ref{lemma-covering-permanence}.

\medskip\noindent
Let us consider condition (3) for $n > 0$. We claim there is an
injective map $\tau : S' \to S$ of finite sets, such that for any
object $T$ of $\text{SR}(\mathcal{C}, X)$ the morphism
\begin{equation}
\label{equation-map}
\Hom(C[k], T)_{n + 1} \to
(\text{cosk}_n \text{sk}_n \Hom(C[k], T))_{n + 1}
\end{equation}
is isomorphic to the projection $\prod_{s \in S} T \to \prod_{s' \in S'} T$
functorially in $T$. If this is true, then we see, arguing as in the previous
paragraph, that the arrow $\gamma$ is the morphism
$$
\prod\nolimits_{s \in S} Z
\longrightarrow
\prod\nolimits_{s \in S'} Z
\times
\prod\nolimits_{s \not\in \tau(S')} Y
$$
which is a product of coverings and hence a covering by
Lemma \ref{lemma-covering-permanence}. By construction, we have
$\Hom(C[k], T)_{n + 1} = \prod_{\alpha : [k] \to [n + 1]} T$
(see Simplicial, Lemma \ref{simplicial-lemma-morphism-into-product}).
Correspondingly we take $S = \text{Map}([k], [n + 1])$.
On the other hand, Simplicial, Lemma \ref{simplicial-lemma-formula-limit},
provides a description of points of
$(\text{cosk}_n \text{sk}_n \Hom(C[k], T))_{n + 1}$
as sequences $(f_0, \ldots, f_{n + 1})$ of points of $\Hom(C[k], T)_n$
satisfying $d^n_{j - 1} f_i = d^n_i f_j$ for $0 \leq i < j \leq n + 1$.
We can write $f_i = (f_{i, \alpha})$ with $f_{i, \alpha}$ a point of $T$
and $\alpha \in \text{Map}([k], [n])$. The conditions translate into
$$
f_{i, \delta^n_{j - 1} \circ \beta} = f_{j, \delta_i^n \circ \beta}
$$
for any $0 \leq i < j \leq n + 1$ and $\beta : [k] \to [n - 1]$. Thus we
see that
$$
S' = \{0, \ldots, n + 1\} \times \text{Map}([k], [n]) / \sim
$$
where the equivalence relation is generated by the equivalences
$$
(i, \delta^n_{j - 1} \circ \beta) \sim (j, \delta_i^n \circ \beta)
$$
for $0 \leq i < j \leq n + 1$ and $\beta : [k] \to [n - 1]$.
A computation (omitted) shows that the morphism (\ref{equation-map})
corresponds to the map $S' \to S$ which sends $(i, \alpha)$ to
$\delta^{n + 1}_i \circ \alpha \in S$. (It may be a comfort to the
reader to see that this map is well defined by part (1) of
Simplicial, Lemma \ref{simplicial-lemma-relations-face-degeneracy}.)
To finish the proof it suffices to show that if
$\alpha, \alpha' : [k] \to [n]$ and $0 \leq i < j \leq n + 1$
are such that
$$
\delta^{n + 1}_i \circ \alpha = \delta^{n + 1}_j \circ \alpha'
$$
then we have $\alpha = \delta^n_{j - 1} \circ \beta$
and $\alpha' = \delta_i^n \circ \beta$ for some $\beta : [k] \to [n - 1]$.
This is easy to see and omitted.
\end{proof}

\begin{proof}[Proof using spectral sequences.]
Suppose that $\xi \in H^p(X, \mathcal{F})$ for some $p \geq 0$.
Let us show that $\xi$ is in the image of the map
$\check{H}^p(X, \mathcal{F}) \to H^p(X, \mathcal{F})$ of
Lemma \ref{lemma-cech-spectral-sequence}
for some hypercovering $K$ of $X$.

\medskip\noindent
This is true if $p = 0$ by Lemma \ref{lemma-h0-cech}.
If $p = 1$, choose a {\v C}ech hypercovering $K$ of $X$ as in
Example \ref{example-cech} starting with a covering
$K_0 = \{U_i \to X\}$ in the site $\mathcal{C}$ such that
$\xi|_{U_i} = 0$, see
Cohomology on Sites,
Lemma \ref{sites-cohomology-lemma-kill-cohomology-class-on-covering}.
It follows immediately from the spectral sequence
in Lemma \ref{lemma-cech-spectral-sequence} that $\xi$ comes
from an element of $\check{H}^1(K, \mathcal{F})$ in this case.
In general, choose any hypercovering $K$ of $X$ such
that $\xi$ maps to zero in $\underline{H}^p(\mathcal{F})(K_0)$
(using Example \ref{example-cech} and
Cohomology on Sites,
Lemma \ref{sites-cohomology-lemma-kill-cohomology-class-on-covering}
again).
By the spectral sequence of Lemma \ref{lemma-cech-spectral-sequence}
the obstruction for $\xi$ to come from an element of
$\check{H}^p(K, \mathcal{F})$ is a sequence of elements
$\xi_1, \ldots, \xi_{p - 1}$ with
$\xi_q \in \check{H}^{p - q}(K, \underline{H}^q(\mathcal{F}))$
(more precisely the images of the $\xi_q$ in certain subquotients
of these groups).

\medskip\noindent
We can inductively replace the hypercovering $K$ by refinements
such that the obstructions $\xi_1, \ldots, \xi_{p - 1}$ restrict to zero
(and not just the images
in the subquotients -- so no subtlety here). Indeed, suppose we have
already managed to reach the situation where
$\xi_{q + 1}, \ldots, \xi_{p - 1}$ are zero.
Note that $\xi_q \in \check{H}^{p - q}(K, \underline{H}^q(\mathcal{F}))$
is the class of some element
$$
\tilde \xi_q \in
\underline{H}^q(\mathcal{F})(K_{p - q}) =
\prod H^q(U_i, \mathcal{F})
$$
if $K_{p - q} = \{U_i \to X\}_{i \in I}$. Let $\xi_{q, i}$
be the component of $\tilde \xi_q$ in $H^q(U_i, \mathcal{F})$.
As $q \geq 1$ we can use
Cohomology on Sites,
Lemma \ref{sites-cohomology-lemma-kill-cohomology-class-on-covering}
yet again to choose coverings $\{U_{i, j} \to U_i\}$
of the site such that each restriction $\xi_{q, i}|_{U_{i, j}} = 0$.
Consider the object $Z = \{U_{i, j} \to X\}$ of the category
$\text{SR}(\mathcal{C}, X)$ and its obvious morphism
$u : Z \to K_{p - q}$. It is clear that $u$ is a covering, see
Definition \ref{definition-covering-SR}. By
Lemma \ref{lemma-covering} there
exists a morphism $L \to K$ of hypercoverings of $X$ such that
$L_{p - q} \to K_{p - q}$ factors through $u$. Then clearly the
image of $\xi_q$ in $\underline{H}^q(\mathcal{F})(L_{p - q})$.
is zero. Since the spectral sequence of
Lemma \ref{lemma-cech-spectral-sequence}
is functorial this means that after replacing $K$ by $L$ we reach the
situation where $\xi_q, \ldots, \xi_{p - 1}$ are all zero.
Continuing like this we end up with a hypercovering where they are all
zero and hence $\xi$ is in the image of the map
$\check{H}^p(X, \mathcal{F}) \to H^p(X, \mathcal{F})$.

\medskip\noindent
Suppose that $K$ is a hypercovering of $X$, that
$\xi \in \check{H}^p(K, \mathcal{F})$ and that the image of
$\xi$ under the map
$\check{H}^p(X, \mathcal{F}) \to H^p(X, \mathcal{F})$ of
Lemma \ref{lemma-cech-spectral-sequence}
is zero. To finish the proof of the theorem we have to show that
there exists a morphism of hypercoverings $L \to K$ such that
$\xi$ restricts to zero in $\check{H}^p(L, \mathcal{F})$.
By the spectral sequence of Lemma \ref{lemma-cech-spectral-sequence}
the vanishing of the image of $\xi$ in $H^p(X, \mathcal{F})$
means that there exist elements $\xi_1, \ldots, \xi_{p - 2}$
with $\xi_q \in \check{H}^{p - 1 - q}(K, \underline{H}^q(\mathcal{F}))$
(more precisely the images of these in certain subquotients)
such that the images $d_{q + 1}^{p - 1 - q, q}\xi_q$ (in the spectral
sequence) add up to $\xi$. Hence by exactly the same mechanism as above
we can find a morphism of hypercoverings $L \to K$ such that
the restrictions of the elements $\xi_q$, $q = 1, \ldots, p - 2$
in $\check{H}^{p - 1 - q}(L, \underline{H}^q(\mathcal{F}))$ are zero.
Then it follows that $\xi$ is zero since the morphism $L \to K$
induces a morphism of spectral sequences according to
Lemma \ref{lemma-cech-spectral-sequence}.
\end{proof}

\bibliographystyle{plain}
\bibliography{my}
\end{document}
